\section{微小扰动为何能翻转预测}
\label{sec:17-Constrained-Guarantees/02-Adversarial-Robustness/01}

一个图像分类器在测试集上做到 \(96\%\)。有人拿这个模型做了一件简单的事：对每张测试图片，把每个像素值改动至多 \(8/255\)——改完之后的图片与原图并排放着，肉眼分不出差别，打印出来完全一样。模型在这批图片上的准确率是 \(4\%\)。

这个结果本身还不算最反直觉的。真正需要解释的是下面这个对照：给同样的图片加上幅度相同的\textbf{随机}噪声——每个像素独立地在 \([-8/255,\,8/255]\) 里均匀抽取——准确率几乎不动，还在九十几。同样的扰动预算，随机地花掉什么也不会发生，按某个方向花掉就能把模型打穿。

所以问题不在``模型对噪声敏感''。模型对噪声一点都不敏感。问题在于扰动空间里存在一些特殊方向，而找到它们并不困难。

\subsection{线性化就能解释大半}

先看一个最简单的情形：二分类线性模型 \(f(x)=w^{\trans}x+b\)。给输入加扰动 \(\delta\)，输出变化恰好是 \(w^{\trans}\delta\)。在 \(\norm{\delta}_\infty\le\varepsilon\) 的约束下，这个量能有多大？

\[
\max_{\norm{\delta}_\infty\le\varepsilon} w^{\trans}\delta
=\varepsilon\sum_{i=1}^{d}\abs{w_i}
=\varepsilon\norm{w}_1,
\qquad
\delta^\star=\varepsilon\,\sign(w).
\]

每个分量各自取到边界、符号与 \(w_i\) 对齐，这就是最优解——\(\ell_\infty\) 约束下线性泛函的最大值由 \(\ell_1\) 范数给出，这两个范数互为对偶，同一个对偶关系在\hyperref[sec:09-Convex-Optimization/07-Applications/02]{``\(L_1\) 正则为何产生稀疏''}里以另一种面貌出现过。

关键在这个式子对维数的依赖。设 \(w\) 的各分量量级相当，都在 \(c\) 附近，那么 \(\norm{w}_1\approx cd\)，扰动能造成的输出变化是 \(\varepsilon c d\)。\(\varepsilon\) 可以小到 \(8/255\approx0.03\)，但 \(d\) 对一张 \(32\times32\) 的彩色图片就是 \(3072\)，对一张常见分辨率的照片是几十万。每个像素只挪动一点点，几十万个``一点点''按同一个方向对齐地累加起来，就是一个大数。

这也解释了随机噪声为什么无效。取 \(\delta\) 的各分量为独立的 \(\pm\varepsilon\) 随机符号，\(w^{\trans}\delta\) 是均值为零的随机变量，标准差是 \(\varepsilon\norm{w}_2\approx\varepsilon c\sqrt d\)。对齐能拿到 \(\varepsilon c d\)，随机只能拿到 \(\varepsilon c\sqrt d\) 的量级，两者相差 \(\sqrt d\) 倍——维数越高，差距越大。而且随机的那个量集中得很好，\hyperref[sec:05-Probability/07-Transforms-and-Bounds/04]{``Chernoff 方法与 Hoeffding 集中界''}给出的尾部估计说明它超出几倍标准差的概率指数小，靠运气撞到有效方向是没希望的。

换个说法：随机方向与 \(w\) 几乎正交。两个高维随机向量的夹角余弦集中在 \(0\) 附近，量级是 \(1/\sqrt d\)，这是高维空间最基本的现象之一，也是\hyperref[sec:03-Linear-Algebra/03-Orthogonality/01]{``正交与正交补：把空间分成互不干扰的方向''}那套语言在概率意义下的版本。对抗扰动做的事情只有一件：不再随机取方向，而是把预算全部投在 \(w\) 张成的那一维上。

\begin{center}
\begin{tikzpicture}[x=1.7cm,y=1.7cm,>=stealth,font=\small]
  \begin{scope}
    \clip (-2.0,-1.3) rectangle (2.5,2.3);
    \fill[blue!7] (-3,-2) rectangle (3,3);
    \fill[red!8] (-3,3.375) -- (3,-1.125) -- (3,3) -- (-3,3) -- cycle;
  \end{scope}
  \draw[very thick,black!70] (-1.565,2.3) -- (2.5,-0.75);
  \node[anchor=west,font=\footnotesize,black!70] at (1.5,-0.38) {决策边界};
  % 训练点
  \foreach \p/\q in {-0.9/-0.6, -1.4/0.5, 0.9/-0.9, -0.35/-0.95}
    \fill[blue!60!black] (\p,\q) circle (1.7pt);
  \foreach \p/\q in {1.6/1.2, 0.4/1.9, 1.95/0.55}
    \fill[red!65!black] (\p,\q) circle (1.7pt);
  % 关注的那个点与它的 ell_infty 球
  \draw[black!60,dashed] (-0.7,-0.7) rectangle (0.7,0.7);
  \fill[black] (0,0) circle (2.2pt);
  \node[anchor=east,font=\footnotesize,black!60] at (-0.85,-0.45)
    {\(\ell_\infty\) 球（半径 \(\varepsilon\)）};
  % 两个扰动
  \draw[->,very thick,red!75!black] (0,0) -- (0.68,0.68);
  \draw[->,very thick,black!55] (0,0) -- (0.68,-0.68);
  \node[anchor=south west,font=\footnotesize,red!75!black] at (0.72,0.74) {对抗方向};
  \node[anchor=north,font=\footnotesize,black!55] at (0.32,-0.80) {随机方向};
\end{tikzpicture}
\end{center}

\emph{要记住的是这张图里方块与边界的关系：两个箭头都止于同一个虚线方块的角上，\(\ell_\infty\) 范数完全相同，一个越过了边界，另一个远远不够。训练点到边界的距离小于方块``对角''方向的伸展，于是在同一个扰动预算内，沿边界法向走能出界，沿别的方向走不能。维数越高，方块的对角伸展相对于边距就越夸张，而图里画不出这一点——二维的方块只比球大一点，高维的方块比球大得离谱。}

\subsection{高维方块比高维球大得离谱}

上一段最后那句话值得算一下，因为它是整件事的几何根源。

\(\ell_\infty\) 球（边长 \(2\varepsilon\) 的方块）里最远的点是顶点，它到中心的 \(\ell_2\) 距离是 \(\varepsilon\sqrt d\)。而 \(\ell_2\) 球的半径就是 \(\varepsilon\)。二维时这个比值是 \(1.41\)，画出来只是方块的角比圆稍微凸出一点；\(d=3072\) 时这个比值是 \(55\)。人们对``肉眼看不出的改动''的直觉是按 \(\ell_\infty\) 建立的（每个像素改一点点），而模型的决策依据是内积，走的是 \(\ell_2\) 与对齐方向的账。两套度量在高维下差了 \(\sqrt d\) 倍，直觉与模型算的根本不是同一个量。

同一个 \(\sqrt d\) 还以另一种方式出现。半径 \(1\) 的高维球，其体积的绝大部分集中在紧贴表面的一层薄壳里：内半径 \(1-t\) 的球体积占比是 \((1-t)^d\)，\(d=3072\)、\(t=0.01\) 时这个数约为 \(5\times10^{-14}\)。也就是说，在这样的维数下``靠近中心''这件事在体积意义上几乎不发生，随机取一个点，它必定贴着边缘。

把这两条放在一起，就得到本节最重要的判断。设数据落在某个半径为 \(R\) 的区域里，训练集有 \(n\) 个点，每个点周围的 \(\varepsilon\)-邻域是一个小球。这些小球的体积之和相对整个区域的占比至多是 \(n(\varepsilon/R)^d\)。取 \(n=10^5\)、\(\varepsilon/R=0.1\)、\(d=3072\)，这个数是 \(10^5\times10^{-3072}\)。训练过程在这样一个体积占比里验证了模型的行为，然后我们指望结论覆盖整个空间。

所以``在训练点上正确''与``在训练点的邻域上正确''不是程度之差，是两个几乎无关的陈述。经验风险最小化只约束了前者——它在训练点处压低损失，对邻域内的行为不作任何要求，这一点在\hyperref[sec:13-Machine-Learning/10-Learning-Theory/01]{``训练误差为何总是乐观''}里已经以另一种形式出现过：训练目标没有要求的事情，优化器不会顺带做到。

\subsection{神经网络不是线性的，但线性化解释仍然管用}

前面的推导用的是线性模型，而被攻破的是深度网络。这一步跳跃需要交代。

在一个输入点附近做一阶展开，\(f(x+\delta)\approx f(x)+\nabla_x f(x)^{\trans}\delta\)，梯度扮演了 \(w\) 的角色。于是最优的一阶扰动方向是 \(\varepsilon\,\sign(\nabla_x f(x))\)，这正是最早也最简单的攻击方法所做的事：一次反向传播拿到输入梯度，取符号，乘以 \(\varepsilon\)。它极其便宜，却已经能把很多模型打到个位数准确率。

一阶展开在 \(\varepsilon\) 很小时才准确，而 \(8/255\) 在像素空间里确实很小。用 ReLU 网络的话还有一层额外的理由：这类网络是分片仿射的，在每个激活模式不变的区域内它\textbf{就是}线性的，一阶展开在该区域内没有误差。扰动幅度小到不跨越太多分片边界时，线性分析给出的就是真实行为。

需要说明这个解释的边界。它解释了``为什么存在有效方向''和``为什么随机方向无效''，这两件事可以算清楚。它没有解释为什么这些方向在很宽的 \(\varepsilon\) 范围内一直有效，也没有解释多步攻击为什么显著强于单步——后者说明一阶展开在实际使用的 \(\varepsilon\) 下已经不够精确。

\begin{remark}
这一节里，哪些是可算的、哪些是观察到的，需要分清。对偶范数那一步、\(\sqrt d\) 的差距、体积占比的估计，都是定理或直接计算，条件明确。而``对抗样本在不同模型之间可迁移''——在模型 A 上构造的扰动往往也能骗过结构不同、训练数据不同的模型 B——是一条被反复复现的\textbf{经验观察}，它是黑盒攻击可行的基础。目前有几种解释（不同模型学到的决策边界在数据流形附近相似、扰动方向对应数据中真实但脆弱的特征），都还没有被公认的证据分出高下。把它当成稳定的经验事实来做安全评估是合理的，把它当成有证明的性质来推理不是。
\end{remark}

同样属于经验规律的还有一条：对抗脆弱性不是训练不充分造成的。更多数据、更长训练、更大模型都不会让它消失，标准训练下的准确率越高，反而往往越容易被攻击。有一种解释认为这是因为模型确实学到了数据中真实存在但幅度很小的判别特征，这些特征泛化得很好却极不稳健，而标准训练没有任何理由不去使用它们。这个解释与\hyperref[sec:14-Deep-Learning/04-Convolutional-Networks/05]{``归纳偏置何时帮忙何时挡路''}里的判断一致——模型使用哪些特征由训练目标决定，目标里没写``只用稳健的特征''，就不要指望它自己挑。

\subsection{于是问题变成了一个全称命题}

把前面几段收拢，鲁棒性要求的是这样一件事：对每个输入 \(x\) 和\textbf{每个}满足 \(\norm{\delta}\le\varepsilon\) 的扰动，预测都不变。这是一个关于邻域的全称命题，而训练只见过有限个点。

这个形状与上一个单元的隐私定义惊人地相似——那里是``对任意相邻数据集、对任意输出集合''，这里是``对每个输入、对邻域内每个扰动''。两处的全称量词都意味着同一件事：验证它不能靠采样。跑一万次攻击都没打穿，不构成``邻域内没有反例''的证据，只说明这一万次没找到。

由此分出两条路。一条是继续找反例，找不到就暂且相信——这是绝大多数``鲁棒性评测''在做的事，它给出的是乐观的估计。另一条是想办法证明邻域内不存在反例，给出一个不依赖攻击方法的半径。下一节走第二条路，会看到有一个量直接控制着这个半径，而它同时也是这条路最大的障碍。
